\documentclass[12pt]{article}
\usepackage[utf8]{inputenc}
\usepackage[T1]{fontenc}
\usepackage{hyperref}
\title{Wemba Wemba Language Resources for K-12 Curriculum Development}
\author{Compiled by Claude Sonnet 5.5 (claude-sonnet-5-5) for Andr\'as Kornai}
\date{}
\begin{document}
\maketitle

\section{Introduction}
Wemba Wemba (also spelled Wamba Wamba; ISO 639-3 xww) is a dormant Australian Aboriginal language of the Victoria/New South Wales Murray River region. No remaining community members use it as a first language. Revival work exists, mainly through the Yarkuwa Indigenous Knowledge Centre in Deniliquin. The sources found are few, and this report lists only what appeared in search results. Sources differ on the spelling of the name and on how the language is classified.

\section{Primary Educational Resources}
\begin{itemize}
\item Deniliquin Indigenous Language Project (Yarkuwa), including the ``Our Languages Our Way'' pilot, which took revitalisation work into local primary and secondary schools: \url{https://www.yarkuwa.org.au/deniliquin-Indigenous-language-project}
\item Wurrekangurr learning lists and printable guides: \url{https://www.yarkuwa.org.au/wurrekangurr}
\item Wurrekangurr resource document: \url{https://yarkuwa.org.au/wurrekangurr/resources/4/document}
\end{itemize}

\section{Linguistic and Reference Materials}
\begin{itemize}
\item Wurrekangurr searchable dictionary (more than 1200 entries imported from Luise A. Hercus's WembaWemba Dictionary, based on interviews from 1962--1970): \url{https://www.yarkuwa.org.au/wurrekangurr}
\item Wikipedia overview of the language: \url{https://en.wikipedia.org/wiki/Wemba_Wemba_language}
\item Omniglot page on Wemba Wemba: \url{https://omniglot.com/writing/wembawemba.htm}
\end{itemize}

\section{Religious and Community Resources}
\begin{itemize}
\item Yarkuwa Indigenous Knowledge Centre, about page: \url{https://www.yarkuwa.org.au/about}
\item Yarkuwa culture site, about page: \url{https://culture.yarkuwa.org.au/about}
\end{itemize}
No scripture translations were found in the search results.

\section{Digital and Audio Resources}
\begin{itemize}
\item Wurrekangurr website, which includes audio of how words are spoken: \url{https://www.yarkuwa.org.au/wurrekangurr}
\item ``Wemba Wemba'' iPad app, introducing the language and made by Swan Hill Primary School students with the Victorian Aboriginal Corporation for Languages: \url{https://apps.apple.com/us/app/id1610933655}
\end{itemize}

\section{Curriculum Development Context}
The language is dormant. The material is a dictionary of about 1200 entries from mid-twentieth-century fieldwork, plus learning lists and a school pilot. A full K-12 curriculum conducted mainly in the language is not possible at present. A realistic course would be a heritage and revitalisation programme built on the Wurrekangurr word lists and audio, led by the Yarkuwa community. Speakers of nearby, related varieties may be relevant, and local community authority over the language should be respected.
\end{document}